\section{Annihilation as an integrated mode limit}
\label{sec:annihilation}
We retain all assumptions and seed classes of Section~\ref{sec:position}.
The obstruction noted there is a point-field obstruction, not yet an
obstruction to convergence after integration. The following construction
resolves the integrated question for every physical input. A stronger
absolute position formula will then be obtained on physical seed images.

\subsection{Finite emitted-mode filters}
Choose finite-support real diagonal contractions on $P_\nu$,
\begin{equation}
 F_M|n\rangle=f_M(n)|n\rangle,\qquad
 0\le f_M(n)\le1,\qquad f_M(n)\longrightarrow1\quad(n\in\mathbb Z).
 \label{eq:mode-filter}
\end{equation}
Set $\mathfrak a_M(x)=\sum_n f_M(n)\psi_n^\nu(x)a_n$.
On each fixed particle sector both this finite sum and its adjoint are
bounded operators, uniformly in $x$ for fixed $M$. The bound may grow with
$M$. In particular, for $\phi\in\mathcal P_N$,
\begin{equation}
 \zeta^-_{M,\phi,u}(y)=\mathfrak a_M^\dagger(x_R)[E_R(u)\phi](y),
 \qquad
 \norm{a_{\zeta^-_{M,\phi,u}}}_2\le C_{M,\phi}(1+|u|)^2.
 \label{eq:ann-test-bound}
\end{equation}
Indeed $\|\mathfrak a_M^\dagger(x)v\|\le
\sqrt{N+1}\big(\sum_n f_M(n)^2|\psi_n^\nu(x)|^2\big)^{1/2}\|v\|$;
the finite-mode envelope and the packet estimate of the previous section
apply. This defines, independently of a physical field operator,
\begin{equation}
 \mathcal I^-_{N,M}(\phi,k)=
 \int_{\mathbb R}g_{\sigma,\tau}(u)\dd u\int_Y
 \ip{\zeta^-_{M,\phi,u}(y)}{[R\mathsf U_{N+1}(u)k](y)}\dd y,
 \quad k\in\cH_{N+1}^{\max,w}(s_*).
 \label{eq:ann-cut-integral}
\end{equation}
Equation~\eqref{eq:position-pairing} proves absolute convergence and
continuity in $k$ at each fixed $M$. It does not supply an $M$-independent
absolute majorant for arbitrary $k$.

On the kinematical right system define
\begin{equation}
 C_{N,M}=\sqrt{N+1}(1_R\otimes P_{\Sym^{N+1}})
 (1_R\otimes F_M\otimes1_{\rm spectators})(T\otimes1).
 \label{eq:filtered-creation}
\end{equation}
The filter acts on the newly emitted leg before Bose symmetrization.
It is not a filter on all output particles. Thus
\begin{equation}
 \norm{C_{N,M}}\le b_N:=\sqrt{(N+1)\mathfrak M},\qquad
 C_{N,M}v\longrightarrow C_Nv\quad(v\in W_N).
 \label{eq:filtered-strong}
\end{equation}
Strong convergence follows from $F_M\to1$ strongly and boundedness of $T$.
Since $F_M$ commutes with rotations, $C_{N,M}$ preserves right total $K$.
It generally fails to intertwine boosts. No physical operator descended
from $C_{N,M}$, and no preservation of physical null seeds at finite $M$,
is assumed.

\subsection{Fixed-left continuity transfers the limit}
On the expanded class $\cE_N^0$ of Section~\ref{sec:fields} write
$\mathcal G_N=J_N\mathcal A_N$. This agrees with $B$ on actual seeds and
satisfies
\begin{equation}
 \norm{\mathcal G_N(|f,0\rangle_L\otimes v)}_{\FP}
 \le B_f\norm v,\qquad v\in W_N^{(0)}.
 \label{eq:expanded-physical-map}
\end{equation}
This bound concerns arbitrary right vectors of total $K=0$; it does not
require any map producing $v$ to intertwine $G$.

\begin{theorem}[Integrated annihilation limit]
For every fixed $\phi\in\mathcal P_N$, the directly defined functionals
\eqref{eq:ann-cut-integral} satisfy
\begin{equation}
 \sup_{\norm k_{\FP}\le1}
 \left|\mathcal I^-_{N,M}(\phi,k)
       -\ip{B\phi}{\mathsf C_N^*k}_{\FP}\right|
 \longrightarrow0.
 \label{eq:ann-dual-limit}
\end{equation}
The limit is independent of every filter sequence obeying
\eqref{eq:mode-filter}. In particular it is the same for sharp symmetric
cutoffs and Fej\'er weights. It is independent of physically null changes
of the test seed, although an individual finite-$M$ functional need not be.
\label{thm:ann-integrated}
\end{theorem}
\begin{proof}
First let $k=B\psi$, $\psi\in\mathcal P_{N+1}$. The joint absolute bound
\eqref{eq:triple-fubini-bound}, now with $\zeta^-_{M,\phi,u}$, is at most
$C_{M,\phi}C_\psi(1+|u|)^4$. Hence the Gaussian permits the interchange
of $u$, $g$, and $y$. At fixed $g$ the emitted sum is finite. Angular
integration selects only finitely many observer modes of the transformed
input; its mass supports remain compact. The original local matrix
element~\eqref{eq:local-mode} therefore gives
\begin{equation}
 \mathcal I^-_{N,M}(\phi,B\psi)
 =\int_G\ip{(1_L\otimes C_{N,M})\phi}{U_{N+1}(g)\psi}_{\rm kin}\dd g.
 \label{eq:ann-finite-ga}
\end{equation}
To check the phase, the time factors on the two sides of the local
pairing give $e^{-iu(m_{\rm out}-m_{\rm in})}$. In
\eqref{eq:ann-finite-ga} the creation coefficient is in the first,
conjugate-linear entry, so its Gaussian phase is conjugated exactly as
required for annihilation. The real Gaussian center is still $\tau$.
No time reversal or change to $-\tau$ has been made.

The first vector in \eqref{eq:ann-finite-ga} lies in $\cE_{N+1}^0$:
the left packet is unchanged and the right total magnetic label is zero.
Absolute scalar group averaging on that class gives
\begin{equation}
 \mathcal I^-_{N,M}(\phi,k)=\ip{\kappa_M(\phi)}k_{\FP},\qquad
 \kappa_M(\phi)=\mathcal G_{N+1}[(1_L\otimes C_{N,M})\phi].
 \label{eq:ann-riesz}
\end{equation}
Initially this identity holds on seed inputs. At fixed $M$, the left side
is continuous in the physical input by \eqref{eq:ann-test-bound}; the
right side is continuous by Hilbert duality. Physical density extends it
to every $k$. In particular we have not commuted $C_{N,M}$ through a
group average.

Write $\phi=\sum_{j=1}^r |f_j,0\rangle_L\otimes v_j$.
For the unfiltered intertwiner the existing descent gives
$\mathcal G_{N+1}[(1_L\otimes C_N)\phi]=\mathsf C_NB\phi$.
Consequently
\begin{equation}
 \norm{\kappa_M(\phi)-\mathsf C_NB\phi}_{\FP}
 \le\sum_{j=1}^r B_{f_j}\norm{(C_{N,M}-C_N)v_j}
 \longrightarrow0.
 \label{eq:ann-error}
\end{equation}
Equations~\eqref{eq:ann-riesz} and \eqref{eq:ann-error} prove
\eqref{eq:ann-dual-limit}. Only the limiting intertwiner is used for
physical descent. The bound in \eqref{eq:ann-error} is a fixed-test bound,
not a uniform bound over all representatives with unit $\norm{B\phi}$.
\end{proof}

\subsection{An absolute uncut formula on seed inputs}
For this stronger statement use specifically
\begin{equation}
 f_M(n)=\left(1-\frac{|n|}{M+1}\right)_+,\qquad
 K_M(\theta)=\frac1{M+1}\left|\sum_{j=0}^M e^{ij\theta}\right|^2.
 \label{eq:fejer-filter}
\end{equation}
The standard positive-kernel argument~\cite[Appendix B]{VasyFourier}
extends to Hilbert-valued continuous functions. For completeness,
$K_M\ge0$, $\int K_M\dd\theta/(2\pi)=1$, and outside any fixed
neighborhood of zero its integral tends to zero, since
$K_M(\theta)\le[(M+1)\sin^2(\theta/2)]^{-1}$.
Uniform continuity and this bound prove uniform convergence of the
convolution to a continuous periodic Hilbert-valued function. The same
positivity bounds the norm of every convolution by its supremum norm.

For a finite occupation-mode vector $v\in\mathcal K_{N+1}$ define
\begin{equation}
 A_{g,v}(x)=V_N(g)\mathfrak a_\nu(g^{-1}x)v.
 \label{eq:covariant-ann-orbit}
\end{equation}
This is a finite sum of smooth transformed scalar modes times fixed
Hilbert vectors, with $\sup_{g,x}\norm{A_{g,v}(x)}\le c_v$.
Covariance of the positive-frequency one-particle representation implies
that its angular Fourier coefficients are those of
$\sum_n\psi_n^\nu(x)a_nV_{N+1}(g)v$. This equality can first be read
after smearing in $x$, where the annihilation operator is defined on
every fixed particle sector; smoothness of \eqref{eq:covariant-ann-orbit}
then identifies the Fourier coefficients. No point evaluation on an
arbitrary one-particle Hilbert vector is needed. Fej\'er convolution gives
\begin{equation}
 \mathfrak a_M(x)V_{N+1}(g)v\longrightarrow A_{g,v}(x),\qquad
 \norm{\mathfrak a_M(x)V_{N+1}(g)v}\le c_v.
 \label{eq:ann-fejer-domination}
\end{equation}
Convergence is uniform in the angle for fixed $g$ and time coordinate;
no uniform convergence in $g$ is claimed.

For a product seed $\psi(y)=u_L(x_L)u_R(x_R)v$ put
$u_{R,u}=E_R(u)u_R$ and define
\begin{equation}
 D_\psi(u;y)=\int_G u_L(g^{-1}x_L)u_{R,u}(g^{-1}x_R)
                    A_{g,v}(x_R)\dd g.
 \label{eq:uncut-seed-ann}
\end{equation}
Extend this definition by linearity to $\mathcal P_{N+1}$.
The scalar absolute orbit integral of the observer product is finite
almost everywhere, by the seed orbit bound. Thus
\eqref{eq:uncut-seed-ann} is a Bochner integral almost everywhere in
$(u,y)$. For finite $M$, boundedness of $\mathfrak a_M(x_R)$ lets it
pass through the seed's Bochner group average. Equation
\eqref{eq:ann-fejer-domination} and dominated convergence then show
\begin{equation}
 D_\psi(u;y)=\lim_{M\to\infty}
 \mathfrak a_M(x_R)[R\mathsf U_{N+1}(u)B\psi](y)
 \quad\text{almost everywhere}.
 \label{eq:point-fejer-limit}
\end{equation}
Joint representatives can be chosen by the local $L^1$ construction in
Section~\ref{sec:position}. Coarea and the two seed orbit bounds give
\begin{equation}
 \int_{\mathbb R}g_{\sigma,\tau}(u)\dd u\int_Y
 \norm{[E_R(u)\phi](y)}\norm{D_\psi(u;y)}\dd y<\infty.
 \label{eq:uncut-ann-absolute}
\end{equation}
Indeed, before the $g$ integral, the upper bound is a constant times
the product of the observer seed envelopes. Its $g,y$ integral is at
most $C_{\phi,\psi}(1+|u|)^4$. This same majorant is independent of $M$
in \eqref{eq:ann-fejer-domination}. We may therefore take the limit
inside the direct integral and apply Theorem~\ref{thm:ann-integrated}:
\begin{equation}
 \int_{\mathbb R}g_{\sigma,\tau}(u)\dd u\int_Y
 \ip{[E_R(u)\phi](y)}{D_\psi(u;y)}\dd y
 =\ip{B\phi}{\mathsf C_N^*B\psi}_{\FP}.
 \label{eq:uncut-ann-identified}
\end{equation}
The pointwise Fej\'er limit depends only on $B\psi$: if $B\psi=0$,
every finite-$M$ expression in \eqref{eq:point-fejer-limit} vanishes
almost everywhere, and hence so does its limit.

\subsection{Scope of the two conclusions}
For arbitrary physical inputs the proved statement is convergence of
integrated functionals in physical dual norm at each fixed seed test.
For the dense physical seed images, \eqref{eq:uncut-ann-absolute} gives
the stronger absolute uncut formula with the pointwise Fej\'er meaning
\eqref{eq:point-fejer-limit}. Neither assertion proves pointwise
convergence or an absolute norm majorant for every physical input.
Sharp cutoffs share the integrated limit but have not been assigned the
positive majorant used for the seed formula. These distinctions complete
the specified weak position interpretation of both field blocks while
leaving unrestricted point multiplication, zero threshold, and removal
of Gaussian smearing open.
